\section{The interval-Krawczyk certification method}
\label{sec:certification}

For each fixed graph $G$ we produce a \emph{replayable} certificate of
the existence of an $\Sphere$-flow on $G$. The certificate is the
output of the four-stage pipeline summarised in this section:
numerical witness search (\cref{sec:certification-witness}), gauge fix
and Newton refinement (\cref{sec:certification-refine}), interval
Krawczyk contraction (\cref{sec:certification-krawczyk}), and a
schema-v2 JSON document (\cref{sec:certification-replay}) that an
independent verifier can re-run from scratch. The mathematical
content sits in
\cref{sec:certification-krawczyk}: the Krawczyk--Moore theorem, when
its contraction hypothesis is satisfied, simultaneously proves
existence and uniqueness of a real zero of the pinned polynomial
system in an explicit box, and therefore certifies an
$\Sphere$-flow on $G$. Implementation pointers: \path{scripts/witness.py},
\path{scripts/exact.py}, \path{scripts/interval.py},
\path{scripts/certificate.py}, \path{scripts/sweep.py},
\path{scripts/verify_sweep.py}.

\subsection{Numerical witness search}
\label{sec:certification-witness}

The first stage produces a \emph{floating-point} approximation to a
zero of the pinned polynomial system $F$. We minimise
$\lVert F(x) \rVert^{2}$ by Levenberg--Marquardt
(\texttt{scipy.optimize.least\_squares} with \texttt{method="lm"})
starting from a random unit-vector assignment on every unpinned edge.
The residual at the returned point gates the verdict: accept as a
witness if $\lVert F \rVert \le 10^{-10}$, otherwise retry from a new
seed. See \path{scripts/witness.py:find_witness} (lines~75--127) for
the multi-seed schedule.

Two failure modes occur in practice:
\begin{enumerate}[label=(\roman*)]
  \item \emph{Saddle plateau.} The LM iterate stalls at a
        non-minimal stationary point of $\lVert F \rVert^{2}$.
        A fresh random seed almost always escapes the plateau.
  \item \emph{Near-singular Jacobian.} On large snarks the
        Jacobian at the iterate is poorly conditioned, the LM step
        becomes numerically uninformative, and the residual
        stagnates above $10^{-10}$. Multi-seed retry with seed
        rotation is the implemented remedy.
\end{enumerate}
The numerical witness is used only as an initial guess for the
Newton refinement of \cref{sec:certification-refine}; its inaccuracy
is irrelevant once the refined centre is fed to the Krawczyk operator,
which only requires the centre to lie inside the box of radius
$10^{-5}$ around the true zero.

\subsection{Gauge fix and Newton refinement}
\label{sec:certification-refine}

We rotate the numerical witness into the pinning gauge of
\cref{sec:prelim-gauge} so that the three spokes at vertex $0$ agree
with the canonical $120$-degree triple, then refine the rotated
witness by Newton's method in \texttt{mpmath} at $50$ decimal digits.
Newton refinement solves
$J(c)\,\Delta = F(c)$ for the update $\Delta$ using
\texttt{mpmath.lu\_solve} at $50$~dps, applied iteratively until
$\lVert F(c) \rVert$ stops decreasing (see
\path{scripts/interval.py:_newton_refine}, lines~144--168). All
rational coefficients are handled exactly inside the ring
$\Q[s]/(s^{2} - 3)$, so the formal symbol $s$ represents $\sqrt{3}$
without round-off until the Newton solve itself. Dropping the
Kirchhoff equation at the last vertex (\cref{lem:square-system})
yields the square system to which Newton is applied.

The empirical residual at the refined centre across all $3247$
certificates is $\lvert F(c) \rvert \approx 10^{-51}$, one full order
of magnitude past the $50$-dps precision floor and ten orders below
the box radius $10^{-5}$ used in \cref{sec:certification-krawczyk}.

\subsection{Interval Krawczyk}
\label{sec:certification-krawczyk}

Let $F : \R^{N} \to \R^{N}$ denote the square pinned polynomial
system of \cref{lem:square-system}, viewed as a map between
$N = 9|V|/2 - 8$-dimensional real spaces, and let $J = J_F$ denote
its Jacobian. Let $c \in \R^{N}$ be the Newton-refined centre of
\cref{sec:certification-refine}, and let
$\mathbf{I} = [c - r, c + r]$ denote the closed box of radius
$r = 10^{-5}$ around $c$ (componentwise). The \emph{Krawczyk operator}
is
\begin{equation}
\label{eq:krawczyk-operator}
   K(\mathbf{I}) \;=\; c \;-\; Y\,F(c) \;+\; (E - Y\,J(\mathbf{I}))\,(\mathbf{I} - c),
\end{equation}
where $E$ is the $N \times N$ identity matrix and $Y$ is any real
matrix close to $J(c)^{-1}$; in the implementation
$Y = J(c)^{-1}$ is computed at $50$~dps. The product $Y\,J(\mathbf{I})$
and the term $J(\mathbf{I})$ are interval matrices obtained by the
natural interval extension of $J$ over $\mathbf{I}$, evaluated in
\texttt{mpmath.iv} arithmetic; the term $F(c)$ is evaluated at the
point $c$ (\emph{not} over $\mathbf{I}$) and then wrapped as a
degenerate interval. See \path{scripts/interval.py:_krawczyk_inclusion}
(lines~171--242).

\begin{theorem}[Krawczyk--Moore]
\label{thm:krawczyk}
\TODO{Add \texttt{refs.bib} entries for Krawczyk 1969 (\emph{Computing}~4,
187--201) and/or Moore 1977; THEOREM.md already cites Krawczyk 1969 in
its References section. Citation should appear here.}
Let $F : \R^{N} \to \R^{N}$ be $C^{1}$ on a neighbourhood of
$\mathbf{I}$. If
$K(\mathbf{I}) \subset \mathrm{int}(\mathbf{I})$ componentwise, then:
\begin{enumerate}[label=(\roman*)]
  \item $F$ has a zero in $\mathbf{I}$;
  \item the zero is unique in $\mathbf{I}$; and
  \item every matrix in $J(\mathbf{I})$ is nonsingular.
\end{enumerate}
\end{theorem}

The implementation tests strict componentwise inclusion: for every
coordinate $i$, $K(\mathbf{I})_i.\text{lo} > c_i - r$ and
$K(\mathbf{I})_i.\text{hi} < c_i + r$. The minimum over $i$ of the
two slack quantities is recorded as the \emph{Krawczyk margin}; the
certificate is \emph{certified} if and only if this margin is
strictly positive (\path{scripts/interval.py:_krawczyk_inclusion},
lines~228--233).

\begin{corollary}[$\Sphere$-flow from a certified zero]
\label{cor:s2-flow-from-zero}
A real zero of $F$ in $\mathbf{I}$ corresponds, by unwinding the
pinning gauge of \cref{sec:prelim-gauge}, to an $\Sphere$-flow on $G$.
Hence if \cref{thm:krawczyk} applies to the pinned system on $G$,
then $G$ admits an $\Sphere$-flow.
\end{corollary}

\paragraph{Numerical gotchas.}
Two implementation points are worth flagging.
\begin{enumerate}[label=(\roman*)]
  \item \emph{Point evaluation of $F(c)$.} If $F(c)$ in
        \cref{eq:krawczyk-operator} is replaced by the interval
        evaluation $F(\mathbf{I})$, the Krawczyk margin grows
        linearly in the box radius and the contraction test fails on
        all but the smallest graphs. The centre term must therefore
        be evaluated at the \emph{point} $c$ in \texttt{mpmath} and
        then wrapped as a degenerate interval; see
        \path{scripts/interval.py} line~204.
  \item \emph{lambdify modules.} The pinned coefficients
        (e.g.\ $\pm\tfrac{1}{2}$) appear as rational numbers in the
        sympy expression for $F$. Calling
        \texttt{sympy.lambdify(\ldots, modules=None)} keeps these as
        Python floats, which represent halves exactly. The
        alternative \texttt{modules='mpmath'} introduces an
        \texttt{mpf} / \texttt{iv.mpf} conversion at every coefficient
        and re-injects rounding error. The certifier uses
        \texttt{modules=None}; see \path{scripts/interval.py}.
\end{enumerate}

\subsection{Schema-v2 replayable certificate}
\label{sec:certification-replay}

For every graph in the catalogue, the certifier emits a JSON document
at schema version $2$ containing exactly the data needed for an
independent replay. The fields are:
\begin{itemize}
  \item the canonical \texttt{graph6} string of $G$;
  \item the canonical edge order used by the orientation convention;
  \item the pinned vertex, the pinned edges, and the canonical
        $120$-degree triple values (with $s$ kept as a formal symbol
        and the relation $s^{2} - 3 = 0$ recorded);
  \item the index of the dropped redundant vertex
        (\cref{lem:square-system});
  \item the list of free variables (in canonical order);
  \item the SHA-256 hash of the polynomial system $F$, computed as a
        canonical string of sympy expressions;
  \item the Newton-refined centre $c$ as decimal strings at $50$~dps;
  \item the box radius $r = 10^{-5}$, per-coordinate Krawczyk
        margins, and the overall minimum margin; and
  \item full provenance metadata: Python, mpmath, sympy, networkx
        versions, OS and hostname, the timestamp, and the git commit.
\end{itemize}
The replay verifier
(\path{scripts/verify_sweep.py:verify_one}, calling
\path{scripts/interval.py:replay_krawczyk}, lines~375--434):
\begin{enumerate}[label=(\arabic*)]
  \item parses the \texttt{graph6} string and rebuilds $G$;
  \item re-runs the orientation and the square-system construction
        from first principles, and checks agreement with the
        certificate's edge order, pinning data, dropped vertex, and
        free-variable list;
  \item recomputes the SHA-256 of the freshly built polynomial
        system and compares to the certificate's stored hash;
  \item re-runs the Krawczyk operator on the parsed centre and box
        radius; and
  \item compares the verdict-of-the-replay against the
        verdict-of-the-certificate.
\end{enumerate}
A certificate replays successfully iff all five steps agree. The
verifier never reads the stored Krawczyk box, margins, or verdict
when computing its own decision; the comparison in step~(5) is the
only point at which the stored verdict is consulted, and only as the
target to match. The cold-environment log
\path{docs/external_replay/} records replay of all $3247$ committed
certificates on an unrelated machine.

A defensive option \texttt{--allow-hash-mismatch} (consumed by
\path{scripts/verify_sweep.py}) relaxes step~(3) to a warning, to
allow cross-version replay when an upstream \texttt{sympy} version
bump changes the canonical string form of the polynomial system
without changing the mathematics.
